% GENERATED FILE. Edit canonical lessons, then run npm run generate:books.
% canonical-source-hash: fnv1a64:0b213d2bcf214958
% canonical-lessons: FR-C348-a-cause-de, FR-C348-grace-a, FR-C348-en-effet, FR-C348-des-lors, FR-C348-origine

\chapter{Grâce à toi}
\label{ch:fr-grace-a-toi}
\hlchaptermodality{348}

\begin{chapteropening}
\textbf{By the end of this chapter:} \emph{I can say what something is due to: because of and thanks to, indeed and from then on, and name its origin.}

Complete the last lesson of chapter 348: I can say what something is due to: because of and thanks to, indeed and from then on, and name its origin.
\end{chapteropening}

\section[à cause de]{à cause de — because of}

\label{lesson:FR-C348-a-cause-de}



\begin{quote}
Before the new one: say the French for consequently, then the French for that is why.
\end{quote}

\subsection*{You'll want to know: à cause de}
\textbf{à cause de} — \textquotedblleft{}because of\textquotedblright{}.

\begin{quote}
Nous sommes restés à la maison à cause de la neige.
\end{quote}

\begin{quote}
\emph{We stayed at home because of the snow.}
\end{quote}

Followed by a noun; with a verb, use \emph{parce que}. \emph{À cause du}, \emph{à cause des}.

\begin{cousinweb}
\emph{à cause de} is built on \emph{cause}, which comes from Latin \emph{causa} \textquotedblleft{}cause, reason, lawsuit\textquotedblright{} (EtymDB; Lewis \& Short).
\end{cousinweb}

\begin{grammarlens}[title={What you've built}]
One more word for telling, explaining and describing at B1.
\end{grammarlens}

\subsection*{Your turn}
\begin{itemize}
\raggedright
  \item \emph{Say it:} \emph{à cause de}
  \item \emph{Say it:} \emph{à cause de}, once more
  \item \emph{Say it:} a sentence of your own with \emph{à cause de}
  \item \emph{Recall:} say the French for consequently, then the French for that is why, then say \emph{à cause de} again
\end{itemize}

\begin{tcolorbox}[breakable,colback=teal!4,colframe=teal!35!black,title={Before you move on}]
What does à cause de mean? (\textquotedblleft{}Because of\textquotedblright{}.) Say it once more.
\end{tcolorbox}
\section[grâce à]{grâce à — thanks to}

\label{lesson:FR-C348-grace-a}



\begin{quote}
Before the new one: say the French for that is why, then the French for because of.
\end{quote}

\subsection*{You'll want to know: grâce à}
\textbf{grâce à} — \textquotedblleft{}thanks to\textquotedblright{}.

\begin{quote}
Grâce à toi, j'ai trouvé un travail.
\end{quote}

\begin{quote}
\emph{Thanks to you, I found a job.}
\end{quote}

For good results; for bad ones, \emph{à cause de}.

\begin{cousinweb}
\emph{grâce} comes through Old French \emph{grace} from Latin \emph{grātia} \textquotedblleft{}favour, thanks\textquotedblright{} (EtymDB; Lewis \& Short), built on \emph{grātus} \textquotedblleft{}pleasing\textquotedblright{}.
\end{cousinweb}

\begin{grammarlens}[title={What you've built}]
One more word for telling, explaining and describing at B1.
\end{grammarlens}

\subsection*{Your turn}
\begin{itemize}
\raggedright
  \item \emph{Say it:} \emph{grâce à}
  \item \emph{Say it:} \emph{grâce à}, once more
  \item \emph{Say it:} a sentence of your own with \emph{grâce à}
  \item \emph{Recall:} say the French for that is why, then the French for because of, then say \emph{grâce à} again
\end{itemize}

\begin{tcolorbox}[breakable,colback=teal!4,colframe=teal!35!black,title={Before you move on}]
What does grâce à mean? (\textquotedblleft{}Thanks to\textquotedblright{}.) Say it once more.
\end{tcolorbox}
\section[en effet]{en effet — indeed; and that is because}

\label{lesson:FR-C348-en-effet}



\begin{quote}
Before the new one: say the French for because of, then the French for thanks to.
\end{quote}

\subsection*{You'll want to know: en effet}
\textbf{en effet} — \textquotedblleft{}indeed; and that is because\textquotedblright{}.

\begin{quote}
Il est fatigué ; en effet, il a travaillé toute la nuit.
\end{quote}

\begin{quote}
\emph{He is tired; indeed, he worked all night.}
\end{quote}

It confirms or explains what came before; it does not mean \emph{in effect}.

\begin{cousinweb}
\emph{effet} comes through Old French \emph{effet} from Latin \emph{effectus} \textquotedblleft{}an effect, a result\textquotedblright{} (EtymDB), from \emph{efficere} \textquotedblleft{}to bring about\textquotedblright{}, built on \emph{facere} \textquotedblleft{}to do, to make\textquotedblright{} (Lewis \& Short).
\end{cousinweb}

\begin{grammarlens}[title={What you've built}]
One more word for telling, explaining and describing at B1.
\end{grammarlens}

\subsection*{Your turn}
\begin{itemize}
\raggedright
  \item \emph{Say it:} \emph{en effet}
  \item \emph{Say it:} \emph{en effet}, once more
  \item \emph{Say it:} a sentence of your own with \emph{en effet}
  \item \emph{Recall:} say the French for because of, then the French for thanks to, then say \emph{en effet} again
\end{itemize}

\begin{tcolorbox}[breakable,colback=teal!4,colframe=teal!35!black,title={Before you move on}]
What does en effet mean? (\textquotedblleft{}Indeed; and that is because\textquotedblright{}.) Say it once more.
\end{tcolorbox}
\section[dès lors]{dès lors — from then on; consequently}

\label{lesson:FR-C348-des-lors}



\begin{quote}
Before the new one: say the French for thanks to, then the French for indeed.
\end{quote}

\subsection*{You'll want to know: dès lors}
\textbf{dès lors} — \textquotedblleft{}from then on; consequently\textquotedblright{}.

\begin{quote}
Il a perdu son travail ; dès lors, tout a changé.
\end{quote}

\begin{quote}
\emph{He lost his job; from then on, everything changed.}
\end{quote}

Two senses: from that moment on, and so. \emph{Dès lors que} means given that.

\begin{cousinweb}
\emph{lors} goes back through Old French to Vulgar Latin \emph{illā hōrā} \textquotedblleft{}at that hour\textquotedblright{} (EtymDB). \emph{dès} \textquotedblleft{}from\textquotedblright{} is usually traced to Latin \emph{dē ex} \textquotedblleft{}from out of\textquotedblright{}.
\end{cousinweb}

\begin{grammarlens}[title={What you've built}]
One more word for telling, explaining and describing at B1.
\end{grammarlens}

\subsection*{Your turn}
\begin{itemize}
\raggedright
  \item \emph{Say it:} \emph{dès lors}
  \item \emph{Say it:} \emph{dès lors}, once more
  \item \emph{Say it:} a sentence of your own with \emph{dès lors}
  \item \emph{Recall:} say the French for thanks to, then the French for indeed, then say \emph{dès lors} again
\end{itemize}

\begin{tcolorbox}[breakable,colback=teal!4,colframe=teal!35!black,title={Before you move on}]
What does dès lors mean? (\textquotedblleft{}From then on; consequently\textquotedblright{}.) Say it once more.
\end{tcolorbox}
\section[l'origine]{l'origine — the origin}

\label{lesson:FR-C348-origine}



\begin{quote}
Before the new one: say the French for indeed, then the French for from then on.
\end{quote}

\subsection*{You'll want to know: l'origine}
\textbf{l'origine} — \textquotedblleft{}the origin\textquotedblright{}.

\begin{quote}
Quelle est l'origine de ce problème ?
\end{quote}

\begin{quote}
\emph{What is the origin of this problem?}
\end{quote}

\emph{À l'origine} — originally. Feminine: \emph{une origine}.

\begin{cousinweb}
\emph{origine} is a learned borrowing from Latin \emph{orīgō} \textquotedblleft{}beginning, source, origin\textquotedblright{} (EtymDB; Lewis \& Short), from \emph{orīrī} \textquotedblleft{}to rise\textquotedblright{}.
\end{cousinweb}

\begin{grammarlens}[title={What you've built}]
One more word for telling, explaining and describing at B1.
\end{grammarlens}

\subsection*{Your turn}
\begin{itemize}
\raggedright
  \item \emph{Say it:} \emph{l'origine}
  \item \emph{Say it:} \emph{l'origine}, once more
  \item \emph{Say it:} a sentence of your own with \emph{l'origine}
  \item \emph{Recall:} say the French for indeed, then the French for from then on, then say \emph{l'origine} again
\end{itemize}

\begin{tcolorbox}[breakable,colback=teal!4,colframe=teal!35!black,title={Before you move on}]
What does l'origine mean? (\textquotedblleft{}The origin\textquotedblright{}.) Say it once more.
\end{tcolorbox}
